\documentclass[12pt]{article}
\usepackage[margin=1in]{geometry}
\usepackage{amsmath, amssymb, amsfonts}
\usepackage{enumitem}
\setlist[enumerate]{leftmargin=*}

\title{Comparative Statics: System, Jacobian, Determinants, and Sign Test}
\author{ECON 320 -- Introduction to Mathematical Economics}
\date{Fall 2025}

\begin{document}
\maketitle

\section*{System and Base Point}
Consider the nonlinear system with parameters $(\alpha,\beta,\gamma)$:
\[
\begin{cases}
F_1:\ x^2 + y + z = \alpha,\\
F_2:\ xy - z = \beta,\\
F_3:\ y^2 + z^2 = \gamma.
\end{cases}
\]
At $(x,y,z)=(1,1,0)$ the system holds for $(\alpha,\beta,\gamma)=(2,1,1)$. We study local effects around this point.

\section*{Step 1. Jacobian with respect to $(x,y,z)$}
\[
J(x,y,z)=
\begin{bmatrix}
\partial_x F_1 & \partial_y F_1 & \partial_z F_1\\
\partial_x F_2 & \partial_y F_2 & \partial_z F_2\\
\partial_x F_3 & \partial_y F_3 & \partial_z F_3
\end{bmatrix}
=
\begin{bmatrix}
2x & 1 & 1\\
y  & x & -1\\
0  & 2y & 2z
\end{bmatrix}.
\]
At $(1,1,0)$,
\[
J=
\begin{bmatrix}
2 & 1 & 1\\
1 & 1 & -1\\
0 & 2 & 0
\end{bmatrix}.
\]

\section*{Step 2. Determinant of the Jacobian and IFT}
Compute $\Delta=\det J$ at $(1,1,0)$:
\[
\Delta
=2\det\!\begin{bmatrix}1&-1\\ 2&0\end{bmatrix}
-1\det\!\begin{bmatrix}1&1\\ 2&0\end{bmatrix}
+0\cdot(\cdots)
=2(2)-(-2)=6.
\]
Since $\Delta=6\neq 0$, the Implicit Function Theorem applies, so $(x,y,z)$ are differentiable functions of $(\alpha,\beta,\gamma)$ locally.

\section*{Step 3. Linearization and Cramer's Rule}
Total differentiation yields
\[
J
\begin{bmatrix}
dx\\ dy\\ dz
\end{bmatrix}
=
\begin{bmatrix}
d\alpha\\ d\beta\\ d\gamma
\end{bmatrix}.
\]
For a single parameter change, set the right side to the corresponding unit vector and use Cramer's rule:
\[
\frac{\partial x}{\partial \theta}=\frac{\Delta_1^{(\theta)}}{\Delta},\quad
\frac{\partial y}{\partial \theta}=\frac{\Delta_2^{(\theta)}}{\Delta},\quad
\frac{\partial z}{\partial \theta}=\frac{\Delta_3^{(\theta)}}{\Delta},\quad
\theta\in\{\alpha,\beta,\gamma\}.
\]

\subsection*{With respect to $\alpha$}
Right-hand side $b_\alpha=(1,0,0)^\top$.
\[
\Delta_1^{(\alpha)}=
\det\!\begin{bmatrix}
1 & 1 & 1\\
0 & 1 & -1\\
0 & 2 & 0
\end{bmatrix}=2,\quad
\Delta_2^{(\alpha)}=
\det\!\begin{bmatrix}
2 & 1 & 1\\
1 & 0 & -1\\
0 & 0 & 0
\end{bmatrix}=0,\quad
\Delta_3^{(\alpha)}=
\det\!\begin{bmatrix}
2 & 1 & 1\\
1 & 1 & 0\\
0 & 2 & 0
\end{bmatrix}=2.
\]
Therefore
\[
\boxed{\ \frac{\partial x}{\partial \alpha}=\frac{1}{3},\quad
\frac{\partial y}{\partial \alpha}=0,\quad
\frac{\partial z}{\partial \alpha}=\frac{1}{3}\ }.
\]

\subsection*{With respect to $\beta$}
Right-hand side $b_\beta=(0,1,0)^\top$.
\[
\Delta_1^{(\beta)}=
\det\!\begin{bmatrix}
0 & 1 & 1\\
1 & 1 & -1\\
0 & 2 & 0
\end{bmatrix}=2,\quad
\Delta_2^{(\beta)}=
\det\!\begin{bmatrix}
2 & 0 & 1\\
1 & 1 & -1\\
0 & 0 & 0
\end{bmatrix}=0,\quad
\Delta_3^{(\beta)}=
\det\!\begin{bmatrix}
2 & 1 & 0\\
1 & 1 & 1\\
0 & 2 & 0
\end{bmatrix}=-4.
\]
Therefore
\[
\boxed{\ \frac{\partial x}{\partial \beta}=\frac{1}{3},\quad
\frac{\partial y}{\partial \beta}=0,\quad
\frac{\partial z}{\partial \beta}=-\frac{2}{3}\ }.
\]

\subsection*{With respect to $\gamma$}
Right-hand side $b_\gamma=(0,0,1)^\top$.
\[
\Delta_1^{(\gamma)}=
\det\!\begin{bmatrix}
0 & 1 & 1\\
0 & 1 & -1\\
1 & 2 & 0
\end{bmatrix}=-2,\quad
\Delta_2^{(\gamma)}=
\det\!\begin{bmatrix}
2 & 0 & 1\\
1 & 0 & -1\\
0 & 1 & 0
\end{bmatrix}=3,\quad
\Delta_3^{(\gamma)}=
\det\!\begin{bmatrix}
2 & 1 & 0\\
1 & 1 & 0\\
0 & 2 & 1
\end{bmatrix}=1.
\]
Therefore
\[
\boxed{\ \frac{\partial x}{\partial \gamma}=-\frac{1}{3},\quad
\frac{\partial y}{\partial \gamma}=\frac{1}{2},\quad
\frac{\partial z}{\partial \gamma}=\frac{1}{6}\ }.
\]

\section*{Step 4. Sign Test Summary}
Since $\Delta>0$, each sign is determined by the numerator's sign:
\[
\begin{array}{lcl}
\alpha:\ (\partial x/\partial\alpha,\ \partial y/\partial\alpha,\ \partial z/\partial\alpha)=(+,\ 0,\ +),\\[4pt]
\beta:\ (\partial x/\partial\beta,\ \partial y/\partial\beta,\ \partial z/\partial\beta)=(+,\ 0,\ -),\\[4pt]
\gamma:\ (\partial x/\partial\gamma,\ \partial y/\partial\gamma,\ \partial z/\partial\gamma)=(-,\ +,\ +).
\end{array}
\]

\end{document}
